\begin{tabularx}{\textwidth}{@{}lX@{}}
\toprule
Term & Meaning in this thesis\\
\midrule
RAG & Retrieval-augmented generation\\
NLP & Natural language processing\\
LLM & Large language model\\
GPT-2 & Tokenizer used for chunk boundaries and lexical evaluation\\
QASPER & Question-answering dataset over scientific research papers\\
OOF & Out-of-fold: a prediction made by a model not fitted on that fold\\
SFT & Supervised fine-tuning\\
MLP & Multilayer perceptron\\
HNSW & Hierarchical navigable small world vector index\\
$q,d$ & Question and its associated source paper\\
$\granularities$ & Candidate chunk sizes $\{10,20,40,80,160\}$ tokens\\
$K$ & Number of retrieved chunks; $K=5$ in the primary protocols\\
$s(q,c)$ & Cosine similarity between a question and a chunk\\
$E(q)$ & Joined reference evidence for a question\\
$R_g(q)$ & Ranked top-$K$ chunks retrieved at size $g$\\
$U(q,g)$ & Joined retrieval token F1 for $R_g(q)$\\
$C(q,g)$ & Regret relative to the best of the five fixed-size actions\\
\bottomrule
\end{tabularx}

\medskip
\enquote{Validation} names the repository split. Because that split was consulted
throughout development, its results are interpreted as development-set
findings, not as a final, untouched test evaluation. \enquote{Classification F1}
and \enquote{retrieval F1} refer to different tasks and are kept separate.
